% g6e_G.tex — 6th Grade Review, Written Explanation Companion.
% Slice: front matter + A1, B1-B6, C1-C7 (source pages 1-11).
% Fragment: \input-ed into a wrapper that already loads mathdocs.sty.

\begin{center}
  {\LARGE\bfseries 6th Grade Review --- Written Explanation Companion}
\end{center}
\medskip

\noindent This companion restates each question from the 6th Grade Review, shows
the checked answer, and explains the work very slowly. Follow the steps in order.
Read every Important words section before you start the numbered steps.

\lesson{A. Number Lines}

\question{A1}

\blockhead{Question}
Plot the point $4$ on the number line below:

\blanknumline{0}{5}{1}

\checked{A solid dot on the tick mark for $4$.}

\blockhead{What the question is asking}
This question asks you to show where the number $4$ sits on a straight line of
numbers by placing a mark there.

\blockhead{Important words}
\begin{words}
  \item A \term{number line}is a straight line with numbers marked in order from left to right.
  \item A \term{tick mark}is a small mark that shows an exact number.
  \item A \term{point}on a number line is one exact place. We show it with a solid filled circle (a dot).
  \item A \term{whole number}is $0, 1, 2, 3, 4$, and so on.
\end{words}

\blockhead{Work the problem one tiny step at a time}
\begin{steps}
  \item Look at the number line printed above. It already has evenly spaced tick marks.
  \item The tick marks are labeled $0$, $1$, $2$, $3$, $4$, $5$ from left to right.
  \item Find the tick mark labeled $4$.
  \item Put a solid filled dot exactly on that tick mark. That means ``the number 4 is here.''
\end{steps}

\numline{0}{5}{1}{4}

\finalans{Solid dot at $4$ on the number line.}

\blockhead{Why this makes sense}
The number $4$ is a whole number, so it sits on a labeled tick, not between
ticks. A common mistake is putting the dot between $3$ and $4$. The correct place
is exactly on the $4$ tick.

\lesson{B. Order Of Operations}

\noindent Before these problems, learn the order-of-operations rule called \textbf{PEMDAS}.

\begin{words}
  \item \textbf{P} means \textit{Parentheses}: do work inside (\;) first.
  \item \textbf{E} means \textit{Exponents}: do powers like $2^{4}$ next.
  \item \textbf{MD} means \textit{Multiply} and \textit{Divide}: do these from left to right.
  \item \textbf{AS} means \textit{Add} and \textit{Subtract}: do these from left to right last.
\end{words}

\question{B1}

\blockhead{Question}
$8 + 3 \times 4 = $\ansblank

\checked{$20$}

\blockhead{What the question is asking}
This question asks for the value of $8 + 3 \times 4$. Do multiply before add.

\blockhead{Important words}
\begin{words}
  \item An \term{expression}is a math phrase with numbers and operations, but no equals sign yet.
  \item \term{Multiply}means ``times,'' written $\times$.
  \item \term{Add}means ``plus,'' written $+$.
  \item \term{PEMDAS}is the order rule: Parentheses, Exponents, Multiply/Divide, Add/Subtract.
\end{words}

\blockhead{Work the problem one tiny step at a time}
\begin{steps}
  \item Look at PEMDAS. There are no parentheses and no exponents in this problem.
  \item Next comes Multiply/Divide. There is a multiply: $3 \times 4$.
  \item Compute the multiply first: $3 \times 4 = 12$.
  \item Now the expression becomes $8 + 12$.
  \item Last comes Add/Subtract. Add: $8 + 12 = 20$.
  \item Check: if you wrongly added first ($8 + 3 = 11$, then $11 \times 4 = 44$),
        that breaks PEMDAS. Multiply comes before add.
\end{steps}

\finalans{$20$}

\blockhead{Why this makes sense}
Multiplication comes before addition in PEMDAS. So we turn $3 \times 4$ into $12$
first, then add $8$. That gives $20$.

\question{B2}

\blockhead{Question}
$(21 - 4) \times 3 = $\ansblank

\checked{$51$}

\blockhead{What the question is asking}
This question asks for the value of $(21 - 4) \times 3$. Do the work inside the
curved marks first, then multiply.

\blockhead{Important words}
\begin{words}
  \item \term{Parentheses}(\;) are grouping symbols. Work inside them first (the P in PEMDAS).
  \item \term{Subtract}means ``take away,'' written $-$.
  \item \term{Multiply}means ``times,'' written $\times$.
\end{words}

\blockhead{Work the problem one tiny step at a time}
\begin{steps}
  \item PEMDAS says P (Parentheses) comes first.
  \item Inside the parentheses: $21 - 4$.
  \item Subtract: $21 - 4 = 17$.
  \item Now the expression is $17 \times 3$.
  \item Multiply: $17 \times 3$. Break it down: $10 \times 3 = 30$ and $7 \times 3 = 21$, then $30 + 21 = 51$.
  \item So $(21 - 4) \times 3 = 51$.
\end{steps}

\finalans{$51$}

\blockhead{Why this makes sense}
The parentheses force you to subtract before you multiply. Without parentheses,
$21 - 4 \times 3$ would be $21 - 12 = 9$. The parentheses change the order, so the
answer is $51$.

\question{B3}

\blockhead{Question}
$2^{4} + 1 = $\ansblank

\checked{$17$}

\blockhead{What the question is asking}
This question asks for the value of $2^{4} + 1$. First build $2^{4}$ by using 2 as
a factor four times, then add 1.

\blockhead{Important words}
\begin{words}
  \item An \term{exponent}(the E in PEMDAS) is a small raised number that tells how
        many times the base is used as a factor.
  \item In $2^{4}$, the \term{base}is 2 and the \term{exponent}is 4. So $2^{4}$ means
        $2 \times 2 \times 2 \times 2$.
  \item A \term{factor}is a number being multiplied. In $2 \times 2 \times 2 \times 2$, each 2 is a factor.
  \item \term{Add}means plus.
\end{words}

\blockhead{Work the problem one tiny step at a time}
\begin{steps}
  \item PEMDAS: no parentheses. Next is E (Exponents).
  \item Compute $2^{4}$: start with 2.
  \item $2 \times 2 = 4$ (that is two factors of 2).
  \item $4 \times 2 = 8$ (that is three factors of 2).
  \item $8 \times 2 = 16$ (that is four factors of 2). So $2^{4} = 16$.
  \item Now the expression is $16 + 1$.
  \item Add: $16 + 1 = 17$.
\end{steps}

\finalans{$17$}

\blockhead{Why this makes sense}
The exponent must be done before the addition. A common mistake is reading $2^{4}$
as $2 \times 4 = 8$. That is wrong. $2^{4}$ means four twos multiplied:
$2 \times 2 \times 2 \times 2 = 16$, then $16 + 1 = 17$.

\question{B4}

\blockhead{Question}
$(5 + 1)^{2} = $\ansblank

\checked{$36$}

\blockhead{What the question is asking}
This question asks for the value of $(5 + 1)^{2}$. Add inside the curved marks
first, then multiply that answer by itself.

\blockhead{Important words}
\begin{words}
  \item \term{Parentheses}come first in PEMDAS.
  \item An \term{exponent}of 2 means square: multiply the number by itself.
  \item So $6^{2}$ means $6 \times 6$.
\end{words}

\blockhead{Work the problem one tiny step at a time}
\begin{steps}
  \item Do P first: inside parentheses, $5 + 1 = 6$.
  \item Now the expression is $6^{2}$.
  \item Do E next: $6^{2} = 6 \times 6$.
  \item Multiply: $6 \times 6 = 36$.
\end{steps}

\finalans{$36$}

\blockhead{Why this makes sense}
Parentheses first, then the square. Do not square 5 and 1 separately.
$(5 + 1)^{2}$ is $6^{2} = 36$, not $5^{2} + 1^{2} = 26$.

\question{B5}

\blockhead{Question}
$3(2^{2}) = $\ansblank

\checked{$12$}

\blockhead{What the question is asking}
This question asks for the value of $3(2^{2})$. A number next to curved marks
means multiply.

\blockhead{Important words}
\begin{words}
  \item $3(2^{2})$ means 3 times the value of $(2^{2})$.
  \item The \term{exponent}2 in $2^{2}$ means $2 \times 2$.
  \item PEMDAS: parentheses first, then exponents inside, then multiply.
\end{words}

\blockhead{Work the problem one tiny step at a time}
\begin{steps}
  \item Look inside the parentheses: $2^{2}$.
  \item Compute the exponent: $2^{2} = 2 \times 2 = 4$.
  \item Now the expression is $3(4)$, which means $3 \times 4$.
  \item Multiply: $3 \times 4 = 12$.
\end{steps}

\finalans{$12$}

\blockhead{Why this makes sense}
The power is inside the parentheses, so square first, then multiply by 3. Notice
this is different from $(3 \times 2)^{2}$, which would square after multiplying.

\question{B6}

\blockhead{Question}
$(3 \times 2)^{2} = $\ansblank

\checked{$36$}

\blockhead{What the question is asking}
This question asks for the value of $(3 \times 2)^{2}$. Multiply inside the curved
marks first, then multiply that answer by itself.

\blockhead{Important words}
\begin{words}
  \item \term{Parentheses}first: multiply inside.
  \item Then the \term{exponent}2 squares that result.
\end{words}

\blockhead{Work the problem one tiny step at a time}
\begin{steps}
  \item Inside parentheses: $3 \times 2 = 6$.
  \item Now the expression is $6^{2}$.
  \item Square: $6^{2} = 6 \times 6 = 36$.
  \item Compare with B5: $3(2^{2}) = 12$, but $(3 \times 2)^{2} = 36$. The
        parentheses change which step happens first.
\end{steps}

\finalans{$36$}

\blockhead{Why this makes sense}
Multiply inside first, then square. A common mix-up is treating this like
$3(2^{2})$. Those are different problems with different answers.

\lesson{C. Decimals, Fractions, And Percents}

\question{C1}

\blockhead{Question}
In $82.731$, the tenths digit is \blank[2.4cm]

\checked{$7$}

\blockhead{What the question is asking}
This question asks which single digit $0$--$9$ sits in the first place after the
dot in $82.731$.

\blockhead{Important words}
\begin{words}
  \item A \term{digit}is a single symbol 0 through 9.
  \item A \term{decimal point}is the dot that separates whole numbers from parts less than one.
  \item \term{Tenths}means the first place to the right of the decimal point (one out of ten equal parts).
  \item \term{Hundredths}is the second place to the right. \term{Thousandths}is the third place to the right.
\end{words}

\blockhead{Work the problem one tiny step at a time}
\begin{steps}
  \item Write the number and name each place: $82.731$.
  \item Left of the point: 8 is tens, 2 is ones.
  \item Right of the point, first digit is tenths: that digit is 7.
  \item Next digit 3 is hundredths. Next digit 1 is thousandths.
  \item The question asks only for tenths, so the answer is 7.
\end{steps}

\finalans{$7$}

\blockhead{Why this makes sense}
Tenths is always the first digit after the decimal point. Do not confuse tenths
($7$) with hundredths ($3$).

\question{C2}

\blockhead{Question}
$314.7 + 29.58 = $\ansblank[1.8cm]\hint{\emph{decimal}}

\checked{$344.28$}

\blockhead{What the question is asking}
This question asks you to add two numbers that use a dot for parts of a whole,
and write the answer the same way.

\blockhead{Important words}
\begin{words}
  \item A \term{decimal}is a number that uses a decimal point to show tenths, hundredths, and so on.
  \item The \term{sum}is the answer to an addition problem.
  \item To \textit{add decimals}, line up the decimal points so equal place values sit in the same column.
  \item To \term{carry}means: when a column adds to 10 or more, write the ones digit
        in that column and move the extra tens digit to the next column on the left.
\end{words}

\blockhead{Work the problem one tiny step at a time}
\begin{steps}
  \item Write the numbers stacked, lining up the decimal points.
  \item Write $314.7$ as $314.70$ so both numbers have two digits after the point.
  \item Add hundredths: $0 + 8 = 8$.
  \item Add tenths: $7 + 5 = 12$. Write 2 in tenths and carry 1 to the ones.
  \item Add ones: $4 + 9 + 1$ (carry) $= 14$. Write 4, carry 1.
  \item Add tens: $1 + 2 + 1$ (carry) $= 4$.
  \item Add hundreds: 3.
  \item Put the decimal point in the sum directly under the other points. The sum is $344.28$:
\end{steps}

\begin{center}
$\begin{array}{r}
  314.70\\
  +\;\,29.58\\ \hline
  344.28
\end{array}$
\end{center}

\finalans{$344.28$}

\blockhead{Why this makes sense}
Lining up decimal points keeps tenths with tenths and hundredths with hundredths.
Check: $314.70 + 29.58$ should be a little more than $314 + 30 = 344$, and
$344.28$ fits.

\question{C3}

\blockhead{Question}
$81.5 - 7.94 = $\ansblank[1.8cm]\hint{\emph{decimal}}

\checked{$73.56$}

\blockhead{What the question is asking}
This question asks you to subtract numbers that use a dot for parts of a whole,
and write the answer the same way.

\blockhead{Important words}
\begin{words}
  \item The \term{difference}is the answer to a subtraction problem.
  \item To \textit{subtract decimals}, line up the decimal points.
  \item You may write extra zeros so both numbers have the same number of decimal places.
  \item To \textit{borrow} (also called \textit{regrouping}) means: when a top digit is
        too small to subtract, take 1 from the next higher place on the left. That
        higher place decreases by 1, and the current place gains 10.
\end{words}

\blockhead{Work the problem one tiny step at a time}
\begin{steps}
  \item Write $81.5$ as $81.50$ to match hundredths.
  \item Line up place values: $81.50$ minus $7.94$. Start digits are tens 8, ones 1, tenths 5, hundredths 0.
  \item Hundredths: top digit 0 is smaller than 4, so borrow from tenths. Rewrite:
        tenths $5 \to 4$ (decrease by 1), hundredths $0 \to 10$ (gain 10). Now $10 - 4 = 6$.
  \item Tenths: after the first borrow we have 4, and 4 is smaller than 9, so borrow
        from ones. Rewrite: ones $1 \to 0$, tenths $4 \to 14$. Now $14 - 9 = 5$.
  \item Ones: after the second borrow we have 0, and 0 is smaller than 7, so borrow
        from tens. Rewrite: tens $8 \to 7$, ones $0 \to 10$. Now $10 - 7 = 3$.
  \item Tens: $7 - 0 = 7$ (no more borrow needed).
  \item Read the rewritten top digits and the difference: $73.56$.
\end{steps}

\begin{center}
\begin{tabular}{r@{\quad}ccccc}
 & tens & ones & . & tenths & hundredths\\ \hline
start top & 8 & 1 & . & 5 & 0\\
after borrows & \textcolor{blue}{7} & \textcolor{blue}{10} & . & \textcolor{blue}{14} & \textcolor{blue}{10}\\ \hline
minus & 0 & 7 & . & 9 & 4\\ \hline
difference & 7 & 3 & . & 5 & 6\\
\end{tabular}

\smallskip
Each blue rewrite shows a borrow: higher place $-1$, current place $+10$.
\end{center}

\finalans{$73.56$}

\blockhead{Why this makes sense}
Check by adding: $73.56 + 7.94 = 81.50$. That matches the starting number, so the
subtraction is correct.

\question{C4}

\blockhead{Question}
$0.32 = $\ansblank[1.8cm]\ \%

\checked{$32\%$}

\blockhead{What the question is asking}
This question asks you to change $0.32$ into an ``out of 100'' amount (written with \%).

\blockhead{Important words}
\begin{words}
  \item A \term{percent}means ``out of 100.'' The symbol is \%.
  \item To change a decimal to a percent, multiply by 100 (move the decimal point
        two places to the right) and write \%.
\end{words}

\blockhead{Work the problem one tiny step at a time}
\begin{steps}
  \item Start with $0.32$.
  \item Multiply by 100: $0.32 \times 100$.
  \item Move the point two places right: $0.32 \to 32$.
  \item Add the percent symbol: $32\%$.
\end{steps}

\finalans{$32\%$}

\blockhead{Why this makes sense}
Percent means hundredths. The decimal $0.32$ is already 32 hundredths, so it is $32\%$.

\question{C5}

\blockhead{Question}
$65\% = $\ansblank[1.8cm]\ as a decimal

\checked{$0.65$}

\blockhead{What the question is asking}
This question asks you to change $65\%$ into a number that uses a dot for parts of a whole.

\blockhead{Important words}
\begin{words}
  \item A \term{percent}is a number out of 100.
  \item A \term{fraction}$\dfrac{a}{b}$ means $a$ equal parts out of $b$ equal parts.
  \item To change a percent to a decimal, divide by 100 (move the decimal point two places to the left).
\end{words}

\blockhead{Work the problem one tiny step at a time}
\begin{steps}
  \item Start with $65\%$. Think of it as $65.0$.
  \item Divide by 100: $65 \div 100$.
  \item Move the point two places left: $65.0 \to 0.65$.
  \item So $65\% = 0.65$.
\end{steps}

\finalans{$0.65$}

\blockhead{Why this makes sense}
$65\%$ means 65 out of 100, which is the fraction $\dfrac{65}{100}$, and that
decimal is $0.65$.

\question{C6}

\blockhead{Question}
$0.8 = $\ansblank[1.8cm]\ as a fraction

\checked{$\dfrac{4}{5}$}

\blockhead{What the question is asking}
This question asks you to write $0.8$ as equal parts over a whole, then divide the
top and bottom by the same common factor until nothing bigger than 1 divides both
evenly.

\blockhead{Important words}
\begin{words}
  \item A \term{fraction}$\dfrac{a}{b}$ means $a$ equal parts out of $b$ equal parts.
  \item The \term{numerator}is the top number. The \term{denominator}is the bottom number.
  \item A \term{factor}of a number divides that number evenly (no leftover).
  \item A \term{common factor}divides both the top and the bottom.
  \item To \term{simplify}(simplifying) a fraction means divide the top and bottom
        by the same common factor to make an equal smaller fraction.
  \item \term{Simplest terms}means the top and bottom share no common factor greater than 1.
  \item The digit in the tenths place means ``over 10.''
\end{words}

\blockhead{Work the problem one tiny step at a time}
\begin{steps}
  \item The decimal $0.8$ has 8 in the tenths place, so $0.8 = \dfrac{8}{10}$.
  \item Find a common factor of 8 and 10. Both divide by 2.
  \item Divide top and bottom by 2: $\dfrac{8 \div 2}{10 \div 2} = \dfrac{4}{5}$.
  \item Check: 4 and 5 share no common factor greater than 1, so $\dfrac{4}{5}$ is simplest.
\end{steps}

\finalans{$\dfrac{4}{5}$}

\blockhead{Why this makes sense}
$0.8$ is eight tenths. Simplifying $\dfrac{8}{10}$ by dividing by 2 gives
$\dfrac{4}{5}$. Check: $4 \div 5 = 0.8$.

\question{C7}

\blockhead{Question}
$15\%$ of $\money{80} = $\ansblank[1.8cm]

\checked{$\money{12}$}

\blockhead{What the question is asking}
This question asks for $15\%$ of $\money{80}$, meaning 15 hundredths of $\money{80}$.

\blockhead{Important words}
\begin{words}
  \item The word \term{of}in percent problems means multiply.
  \item $15\%$ as a decimal is $0.15$ (divide by 100).
\end{words}

\blockhead{Work the problem one tiny step at a time}
\begin{steps}
  \item Change $15\%$ to a decimal: $15 \div 100 = 0.15$.
  \item Multiply: $0.15 \times \money{80}$.
  \item First compute $15 \times 80$. Break it: $15 \times 8 = 120$, then
        $15 \times 80 = 1200$ (one extra zero because $80 = 8 \times 10$).
  \item There are two decimal places in $0.15$, so place the point two places from
        the right in 1200: $12.00$.
  \item So $0.15 \times \money{80} = \money{12}$.
  \item Another way: $10\%$ of $\money{80}$ is $\money{8}$, and $5\%$ of $\money{80}$
        is $\money{4}$, so $15\%$ is $\money{8} + \money{4} = \money{12}$.
\end{steps}

\finalans{$\money{12}$}

\blockhead{Why this makes sense}
Fifteen percent of $\money{80}$ is a little more than one tenth of $\money{80}$. One
tenth of $\money{80}$ is $\money{8}$, so $\money{12}$ is a sensible size.
